\subsection{Experimentos com \textit{Few-Shot Prompting}}

Os experimentos com \textit{few-shot prompting}\footnote{Disponível em: \url{https://github.com/giovannabbottino/prompt-based}.} foram realizados por meio de uma API desenvolvida com o \textit{framework} Flask. O fluxo recebe um texto de entrada, combina-o com \textit{prompts} previamente definidos e encaminha a requisição ao modelo de linguagem executado por meio do Ollama. Como resultado, o serviço disponibiliza o texto originalmente recebido e o grafo de conhecimento validado, serializado em RDF/Turtle.

\begin{figure}[htbp]
    \centering
    \caption{Fluxo de criação de grafos de conhecimento com \textit{few-shot prompting}.}
    \includegraphics[width=0.2\linewidth]{metodo/figuras/experimento-a.jpg}
    \label{fig:experimento-a}
    \source{da autora.}
\end{figure}

Conforme apresentado na Figura~\ref{fig:experimento-a}, o fluxo de criação do grafo de conhecimento inicia-se com o recebimento do texto pela rota \texttt{POST /analyze}. Além do conteúdo textual, a requisição pode incluir parâmetros opcionais relacionados aos arquivos de \textit{prompt} e ao número máximo de tentativas de geração do RDF. Quando esses parâmetros não são informados, são utilizadas as configurações padrão da aplicação.

Em seguida, o serviço carrega o \textit{prompt} de sistema e o \textit{prompt few-shot}. O texto recebido é inserido no \textit{prompt few-shot} por meio da substituição do marcador \texttt{\$\{USER\_TEXT\}}. Após a composição da mensagem, os \textit{prompts} são encaminhados ao Ollama pela rota \texttt{/api/generate}, com o parâmetro \texttt{stream} definido como \texttt{false}, de modo que a resposta seja recebida em um único objeto.

A cada chamada ao modelo, a resposta gerada e os metadados fornecidos pelo Ollama são armazenados em um arquivo CSV. Entre as informações registradas estão o modelo utilizado, a data de criação da resposta, o conteúdo produzido, a duração da geração e as quantidades de \textit{tokens} processados e gerados. Esses registros permitem conservar as informações referentes às diferentes execuções realizadas durante os experimentos.

Após o recebimento da saída do modelo, inicia-se o processo de extração e validação da representação RDF/Turtle. Inicialmente, o serviço identifica o conteúdo correspondente à serialização Turtle e remove, quando presentes, delimitadores de blocos de código e conteúdos textuais anteriores às declarações RDF.

Antes da validação sintática, são produzidos diferentes candidatos de reparo destinados ao tratamento de erros recorrentes nas respostas geradas pelo modelo. Entre os procedimentos implementados estão a normalização de aspas tipográficas, a correção de aspas duplicadas, a inclusão de ponto final quando ausente, a declaração do prefixo \texttt{xsd:} quando utilizado sem definição prévia e o tratamento de identificadores da Wikidata representados de forma inválida. Também podem ser removidos blocos finais incompletos quando já houver declarações RDF completas na resposta.

Na última tentativa de validação, o fluxo também permite uma estratégia de recuperação parcial, na qual são preservadas apenas as declarações que podem ser analisadas sintaticamente. Dessa forma, partes válidas de uma saída parcialmente malformada podem ser mantidas em vez de todo o conteúdo ser descartado.

Cada candidato resultante desses procedimentos é submetido à biblioteca \texttt{rdflib} para análise no formato Turtle. Para que o resultado seja aceito, devem ser atendidas três condições: a representação deve ser sintaticamente válida, conter pelo menos uma tripla e apresentar ao menos uma relação semântica distinta de \texttt{rdfs:label}. Assim, uma saída formada exclusivamente por declarações de rótulos não é considerada um grafo válido para os experimentos.

Quando nenhum candidato produzido em uma tentativa é aceito e ainda existem tentativas disponíveis, uma nova solicitação é encaminhada ao LLM. Nessa chamada, o modelo recebe novamente o \textit{prompt} original, a saída RDF considerada inválida e uma versão limitada da mensagem de erro produzida pelo analisador. A instrução adicional solicita exclusivamente uma versão corrigida do RDF/Turtle, sem explicações ou outros conteúdos textuais.

O parâmetro \texttt{max\_rdf\_attempts} controla o número de gerações destinadas à obtenção de uma representação RDF válida. O valor padrão é três, e a aplicação restringe esse parâmetro ao intervalo de uma a três tentativas. Caso nenhuma tentativa resulte em um grafo válido, o processamento é encerrado e a API retorna o código HTTP \texttt{422}. Quando a validação é concluída com sucesso, a resposta da rota contém os campos \texttt{text}, correspondente ao texto original, e \texttt{rdf}, contendo a serialização Turtle validada.


\subsubsection{Definição dos \textit{prompts}}

Nesse experimento, foram utilizados dois \textit{prompts} com funções complementares: um \textit{prompt} de sistema e um \textit{prompt few-shot}. O primeiro estabelece o papel atribuído ao modelo e as regras gerais para a construção do grafo de conhecimento, enquanto o segundo apresenta as instruções específicas da tarefa, os prefixos autorizados, os exemplos de entrada e saída e o marcador destinado à inserção do texto analisado.

O \textit{prompt} de sistema, apresentado no Apêndice~\ref{apendiceA}, define o modelo como um assistente responsável pela produção de dados estruturados em RDF/Turtle. Entre suas instruções, estabelece-se que entidades reconhecíveis do mundo real devem utilizar identificadores da Wikidata somente quando o respectivo \texttt{Q-ID} for conhecido com confiança. A criação de identificadores inexistentes ou incompletos não é permitida; quando o identificador não é conhecido, devem ser utilizados recursos locais pertencentes ao \textit{namespace} \texttt{kg:}.

O \textit{prompt} também determina que os recursos representados no grafo possuam \texttt{rdfs:label} com marcação de idioma e orienta o modelo a priorizar relações entre entidades. As relações de classificação são preferencialmente representadas pelo predicado \texttt{kg:is}, evitando-se o emprego da abreviação Turtle \texttt{a}. Também são estabelecidas regras relacionadas aos prefixos permitidos, à construção dos predicados, à formação dos identificadores locais e à finalização sintaticamente válida dos blocos de sujeito.

O \textit{prompt few-shot}, apresentado no Apêndice~\ref{apendiceB}, contém as instruções específicas da tarefa de transformação do campo \texttt{Text} em um grafo de conhecimento serializado em RDF/Turtle. A saída solicitada deve conter exclusivamente o RDF, sem notas, títulos, explicações ou delimitadores de blocos de código.

Esse \textit{prompt} complementa as regras definidas pelo \textit{prompt} de sistema ao estabelecer restrições específicas para a serialização, como o encerramento dos blocos de sujeito com ponto final, a utilização de apenas um par de aspas duplas para cada rótulo, a inclusão das marcações de idioma, o emprego de nomes locais em \textit{snake\_case} e a utilização exclusiva de prefixos previamente declarados.

Também são fornecidas instruções relacionadas à representação semântica do conteúdo. O modelo deve considerar as entidades e os diferentes sentidos identificados no texto; nomes alternativos devem ser representados como rótulos adicionais do mesmo recurso; e as relações entre entidades devem ser priorizadas em relação a descrições exclusivamente literais. O \textit{prompt} apresenta ainda um conjunto de predicados preferenciais, como \texttt{kg:is}, \texttt{kg:of}, \texttt{kg:from}, \texttt{kg:in}, \texttt{kg:located\_in} e \texttt{kg:instance\_of}, permitindo o emprego de predicados mais específicos quando derivados do conteúdo analisado.


\subsubsection{Exemplos utilizados no \textit{few-shot prompting}}

A estratégia de \textit{few-shot prompting} foi estruturada mediante a inclusão de exemplos explícitos de transformação de texto em RDF/Turtle. Esses exemplos fornecem ao modelo referências sobre a estrutura esperada das triplas, o uso de identificadores da Wikidata, a construção de recursos locais e a representação das relações semânticas.

Foram utilizados três exemplos. O primeiro representa uma relação entre uma pessoa, um local e um evento histórico; o segundo apresenta uma organização associada ao lançamento de um modelo de inteligência artificial e ao respectivo domínio de atuação; e o terceiro descreve uma descoberta científica envolvendo pessoas, uma substância e um local.

Após os exemplos, o marcador \texttt{\$\{USER\_TEXT\}} delimita a posição destinada ao texto analisado. Durante a execução, esse marcador é substituído pelo conteúdo recebido pela API, enquanto as demais instruções e exemplos permanecem inalterados. Dessa forma, as chamadas utilizam a mesma estrutura de orientação, variando apenas o texto submetido à geração do grafo.


\subsubsection{Configuração do experimento}

A Tabela~\ref{tab:llm-config-a} apresenta os principais parâmetros utilizados no experimento com \textit{few-shot prompting}. São apresentados o nome do parâmetro, seu tipo, sua função e o valor utilizado.

\begin{table}[htbp]
\centering
\caption{Configurações utilizadas no experimento com \textit{few-shot prompting}.}
\label{tab:llm-config-a}
\small
\renewcommand{\arraystretch}{1.3}

\begin{tabular}{
p{3.5cm}
p{1.8cm}
p{6.5cm}
p{2.5cm}
}
\toprule
\textbf{Parâmetro} &
\textbf{Tipo} &
\textbf{Descrição} &
\textbf{Valor utilizado} \\
\midrule

\texttt{stream}
&
\textit{boolean}
&
Desabilita o envio incremental da resposta
&
\texttt{false}
\\

\texttt{seed}
&
\textit{integer}
&
Semente utilizada durante a geração
&
\texttt{42}
\\

\texttt{temperature}
&
\textit{float}
&
Controla o grau de aleatoriedade da geração
&
\texttt{0}
\\

\texttt{top\_k}
&
\textit{integer}
&
Limita a seleção aos \textit{tokens} mais prováveis
&
\texttt{40}
\\

\texttt{top\_p}
&
\textit{float}
&
Define o limiar da amostragem \textit{nucleus}
&
\texttt{0.9}
\\

\texttt{min\_p}
&
\textit{float}
&
Define a probabilidade mínima considerada
&
\texttt{0.05}
\\

\texttt{num\_predict}
&
\textit{integer}
&
Número máximo de \textit{tokens} produzidos pelo modelo em cada geração
&
\texttt{1536}
\\

\texttt{timeout seconds}
&
\textit{integer}
&
Tempo máximo de espera da requisição, em segundos
&
\texttt{180}
\\

\bottomrule
\end{tabular}
\end{table}

Os parâmetros apresentados na Tabela~\ref{tab:llm-config-a} são fornecidos ao cliente Ollama por meio das configurações da aplicação. Os parâmetros opcionais sem valor definido, como \texttt{stop} e \texttt{num\_ctx}, não são incluídos no objeto \texttt{options} enviado na requisição. O parâmetro \texttt{stream}, por sua vez, é definido diretamente como \texttt{false} no corpo da chamada à API do Ollama.

Em síntese, o experimento com \textit{few-shot prompting} utiliza exemplos previamente definidos para orientar a transformação do texto em RDF/Turtle. O LLM é responsável pela geração inicial da representação, enquanto os procedimentos locais realizam a extração, o reparo e a validação sintática do RDF produzido. O grafo somente é retornado pela API após a obtenção de uma representação considerada válida.

